\section{Related Work}
\label{sec:relatedwork}

\textbf{SEP ML Forecasting}
Machine learning has been applied to SEP forecasting mainly as an event-occurrence problem, using solar-source precursors.
\citeA{boubrahimi2017prediction} used other satellites  like GOES and their proton and X-ray fluxes data as multivariate time series for a  interpretable decision trees to predict $>100$ MeV SEP events
\citeA{engell2017sprints} developed SPRINTS, which also uses GOES X-ray flare metadata using statistical, and machine-learned decision tree models.
\citeA{laurenza2018esperta} adapted the ESPERTA empirical model, using flare location, soft X-ray fluence, and 1 MHz radio fluence, providing short-term warnings of moderate-to-extreme solar proton events.
\citeA{aminalragia2021sep} trained neural networks on GOES soft X-ray measurements of solar flares to predict whether a flare would be followed by an SEP event.
\citeA{lavasa2021predictability} treated SEP prediction as a is happening/ is not happening result and compared several machine learning methods on flare and coronal mass ejection data. CME speed,
CME width, and flare soft X-ray fluence were concluded to be informative features. \citeA{kasapis2022interpretable} used support vector machines and regression models on active region magnetograms from SOHO/MDI
to forecast SEP events during Solar Cycle 23. These studies predict whether an event will occur from solar-source precursors, using decision trees, support vector machines,
ensembles, and feedforward neural networks. In contrast, this study forecasts the time profile of the >10 MeV proton flux from relativistic electrons measured in situ at L1
\cite{posner2007forecasting,torres2025machine},
and evaluates attention-based transformer architectures, which have not been applied to this task.

\textbf{Time Series Transformer Applications}
Since their introduction for machine translation \cite{vaswani2017attention}, transformers have been adapted to time series forecasting across many fields \cite{zhao2026transformersurvey}.
Time Series Transformer uses learnable positional encodings with unsupervised pre-training to perform well when data are scarce \cite{zerveas2020transformer}.
More recent models such as Chronos treat forecasting as a language-modeling task and pre-train on many datasets to forecast without task-specific training \cite{ansari2024chronos}.
These architectures have been applied to stock market prediction, clinical and physiological monitoring, battery state-of-health estimation,
and traffic flow forecasting \cite{cai2020traffictransformer,zhao2026transformersurvey}, but rarely to space weather.
Their advantage is not guaranteed, however. Simple baselines can match or beat them: a naive forecast that
repeats the last observed value outperformed transformer models on exchange-rate data \cite{zeng2022transformers},
and transformers tend to overfit when training data are limited \cite{zhao2026transformersurvey}.
\citeA{zhao2026transformersurvey} further note that their performance depends strongly on the choice of positional encoding.